\documentclass[11pt,a4paper]{article}
\usepackage[utf8]{inputenc}
\usepackage[T1]{fontenc}
\usepackage[spanish,es-minimal,es-noshorthands]{babel}
\AtBeginDocument{\renewcommand{\sin}{\operatorname{sen}}\renewcommand{\arcsin}{\operatorname{arcsen}}}
\usepackage{lmodern}
\usepackage[sfdefault,lining]{FiraSans}
\usepackage[a4paper,top=2.3cm,bottom=2.2cm,left=1.8cm,right=1.8cm,headheight=1.3cm,headsep=0.4cm,footskip=1cm]{geometry}
\usepackage{amsmath,amssymb}
\usepackage{icomma}
\usepackage{graphicx}
\usepackage[table]{xcolor}
\usepackage{tikz}
\usepackage{fontawesome5}
\usepackage{ccicons}
\usepackage{fancyhdr}
\usepackage{lastpage}
\usepackage{tasks}
\usepackage{needspace}
\usepackage[most]{tcolorbox}
\usepackage{array}
\usepackage[hidelinks]{hyperref}
\AtBeginDocument{\hypersetup{pdftitle={Ficha de ejercicios de la unidad 4: Ecuaciones. 4º ESO},pdfsubject={Matemáticas A 4º ESO},pdfkeywords={ecuaciones, ecuaciones de primer grado, denominadores, problemas, ecuaciones de segundo grado, completas, incompletas, discriminante, ecuaciones factorizadas, ecuaciones racionales, ecuaciones con radicales}}}

\definecolor{navy}{HTML}{1B2748}
\definecolor{accent}{HTML}{E8603A}
\definecolor{teal}{HTML}{1F9E90}
\definecolor{soft}{HTML}{EEF1F8}
\definecolor{muted}{HTML}{5D6685}

\setlength{\parindent}{0pt}
\setlength{\parskip}{3pt}
\settasks{label=\alph*), label-format=\bfseries\color{navy}, label-width=1.3em, item-indent=1.8em, column-sep=0.6em, after-item-skip=5pt, before-skip=3pt, after-skip=4pt}

% ---------- capçalera i peu ----------
\newcommand{\llicencia}{\href{https://creativecommons.org/licenses/by-nc-sa/4.0/deed.es}{\ccbyncsa}\ \ Licencia Creative Commons Atribución-NoComercial-CompartirIgual 4.0 Internacional (CC BY-NC-SA 4.0)}
\pagestyle{fancy}
\fancyhf{}
\fancyhead[L]{\raisebox{-0.35\height}{\includegraphics[height=0.95cm]{logo-gva.png}}}
\fancyhead[C]{\small\color{muted}Ficha de ejercicios · Unidad 4: Ecuaciones · 4\textsuperscript{o} ESO A}
\fancyhead[R]{\raisebox{-0.35\height}{\includegraphics[height=0.95cm]{logo-vermellar.png}}}
\fancyfoot[L]{\scriptsize\color{muted}\llicencia}
\fancyfoot[R]{\small\color{muted}\thepage}
\renewcommand{\headrulewidth}{0.4pt}
\renewcommand{\headrule}{\color{muted!40}\hrule height \headrulewidth}
\fancypagestyle{primera}{\fancyhf{}\renewcommand{\headrulewidth}{0pt}
  \fancyfoot[L]{\scriptsize\color{muted}\llicencia}
  \fancyfoot[R]{\small\color{muted}\thepage}}

% ---------- nivells i exercicis ----------
\newcommand{\estrelles}[1]{{\color{accent}\ifcase#1\or\faStar\faStar[regular]\faStar[regular]\or\faStar\faStar\faStar[regular]\or\faStar\faStar\faStar\fi}}
\newcounter{exer}
\newcommand{\ex}[1]{\par\needspace{4\baselineskip}\medskip\refstepcounter{exer}\noindent{\color{navy}\bfseries\theexer.}\ \ {\small\estrelles{#1}}\ \ }

\newcommand{\seccio}[2]{\par\needspace{16\baselineskip}\bigskip
  \begin{tcolorbox}[colback=navy,colframe=navy,coltext=white,arc=2mm,boxrule=0pt,left=3mm,right=3mm,top=1.2mm,bottom=1.2mm]
  {\large\bfseries #1\quad #2}
  \end{tcolorbox}\vspace{-2pt}}

\newcommand{\exf}[1]{\par\medskip\refstepcounter{exer}\noindent{\color{navy}\bfseries\theexer.}\ \ {\small\estrelles{#1}}\ \ }
\newcolumntype{C}[1]{>{\centering\arraybackslash}p{#1}}
% quadre de recordatori (només a la fitxa d'exercicis)
\newcommand{\recorda}[1]{\par\smallskip\begin{tcolorbox}[colback=accent!7,colframe=accent!55,arc=2mm,boxrule=0.5pt,left=2.5mm,right=2.5mm,top=1mm,bottom=1mm]\small{\color{accent}\faLightbulb[regular]}\ \textbf{Recuerda:}\ #1\end{tcolorbox}}
\newcommand{\res}[1]{#1}\newcommand{\rest}[1]{#1}

\def\fmneed{11}% que el títol d'una secció no quede sol al final del full
% =====================================================================
% Capçalera FentMates de les fitxes: un color per a cada curs.
% Els colors són els de la pàgina inicial (LEVELS a assets/js/materials.js);
% el logo ve de _eines/logo/pdf/ (python3 _eines/logo/genera.py).
%
% Ús, al final del preàmbul de la unitat:
%   % =====================================================================
% Capçalera FentMates de les fitxes: un color per a cada curs.
% Els colors són els de la pàgina inicial (LEVELS a assets/js/materials.js);
% el logo ve de _eines/logo/pdf/ (python3 _eines/logo/genera.py).
%
% Ús, al final del preàmbul de la unitat:
%   % =====================================================================
% Capçalera FentMates de les fitxes: un color per a cada curs.
% Els colors són els de la pàgina inicial (LEVELS a assets/js/materials.js);
% el logo ve de _eines/logo/pdf/ (python3 _eines/logo/genera.py).
%
% Ús, al final del preàmbul de la unitat:
%   \input{../comu/fentmates.tex}
%   \fmcurs{1eso}
% i, al començament del document:
%   \capcalera{Matemàtiques}{1r ESO}{Fitxa d'exercicis · Unitat 2}{Divisibilitat}
% Opcional abans de \input: \def\fmneed{14} (espai mínim abans de cada secció).
% =====================================================================
\makeatletter
\def\fm@logodir{../../logo/pdf}
\providecommand{\fmneed}{10}
\def\fm@cursnom{}

% clau del curs, color i color del text damunt del color
\newcommand{\fm@curs}[3]{\@namedef{fm@color@#1}{#2}\@namedef{fm@text@#1}{#3}}
\fm@curs{1eso}{E6B84A}{navy}
\fm@curs{2eso}{3B82F6}{white}
\fm@curs{3eso}{FF7A50}{white}
\fm@curs{4eso-a}{A78BFA}{white}
\fm@curs{4eso-b}{6D28D9}{white}
\fm@curs{1batx-mi}{9B1C1C}{white}
\fm@curs{1batx-ccss}{EF6B6B}{white}
\fm@curs{2batx-mii}{1E6B47}{white}
\fm@curs{2batx-ccss}{2BB3A3}{white}

\newcommand{\fmcurs}[1]{%
  \@ifundefined{fm@color@#1}{\PackageError{fentmates}{Curs desconegut: #1}{}}{}%
  \edef\fm@tmp{\noexpand\definecolor{curs}{HTML}{\@nameuse{fm@color@#1}}%
               \noexpand\colorlet{curstext}{\@nameuse{fm@text@#1}}}\fm@tmp
  \colorlet{cursfosc}{curs!72!black}%
  \def\fm@logo{\fm@logodir/vertical-#1.pdf}%
  \def\fm@marca{\fm@logodir/marca-#1.pdf}}

% ---------- capçalera de la primera pàgina ----------
% #1 matèria (petit), #2 curs, #3 tipus de document i unitat, #4 títol
\newcommand{\capcalera}[4]{%
  \gdef\fm@cursnom{#2}%
  \par\noindent
  \begin{tcolorbox}[colback=curs,colframe=curs,boxrule=0pt,arc=2.2mm,width=3.5cm,height=2.15cm,
      valign=center,left=3mm,right=2mm,top=1mm,bottom=1mm,nobeforeafter,box align=center]
    \color{curstext}{\footnotesize #1}\par\vspace{1pt}{\fontsize{19}{21}\selectfont\bfseries #2\par}%
  \end{tcolorbox}\hspace{5mm}%
  \begin{minipage}[c]{\dimexpr\linewidth-3.5cm-5mm-2.9cm\relax}
    {\color{cursfosc}\small\bfseries\MakeUppercase{#3}}\par\vspace{2pt}
    {\color{navy}\fontsize{20}{23}\selectfont\bfseries #4\par}
  \end{minipage}\hfill
  \begin{minipage}[c]{2.6cm}\raggedleft\includegraphics[width=2.6cm]{\fm@logo}\end{minipage}\par
  \vspace{2pt}}

% ---------- barres de secció del color del curs ----------
\renewcommand{\seccio}[2]{\par\needspace{\fmneed\baselineskip}\bigskip
  \begin{tcolorbox}[colback=curs,colframe=curs,coltext=curstext,arc=2mm,boxrule=0pt,left=3mm,right=3mm,top=1.2mm,bottom=1.2mm]
  {\large\bfseries #1\quad #2}
  \end{tcolorbox}\vspace{-2pt}}

% ---------- estrelles de dificultat del color del curs (a Batxillerat no se'n posen) ----------
\@ifundefined{estrelles}{}{%
  \renewcommand{\estrelles}[1]{{\color{curs}\ifcase#1\or\faStar\faStar[regular]\faStar[regular]\or\faStar\faStar\faStar[regular]\or\faStar\faStar\faStar\fi}}}

% ---------- capçalera i peu de les pàgines ----------
% primera pàgina: logos institucionals al peu
\fancypagestyle{primera}{\fancyhf{}\renewcommand{\headrulewidth}{0pt}
  \fancyfoot[L]{\raisebox{-0.32\height}{\includegraphics[height=0.78cm]{logo-gva.png}}\hspace{5mm}%
    \raisebox{-0.32\height}{\includegraphics[height=0.78cm]{logo-vermellar.png}}}
  \fancyfoot[R]{\scriptsize\color{muted}\href{https://creativecommons.org/licenses/by-nc-sa/4.0/deed.ca}{\ccbyncsa}\ \ CC BY-NC-SA 4.0 · fentmates.org\qquad\small\thepage}}
% la resta: el símbol a l'esquerra i el curs a la dreta, amb el títol al centre com abans
\fancyhead[L]{\raisebox{-0.3\height}{\includegraphics[height=0.62cm]{\fm@marca}}}
\fancyhead[R]{\small\bfseries\color{cursfosc}\fm@cursnom}
\renewcommand{\headrule}{\color{curs!45}\hrule height \headrulewidth}
\makeatother

%   \fmcurs{1eso}
% i, al començament del document:
%   \capcalera{Matemàtiques}{1r ESO}{Fitxa d'exercicis · Unitat 2}{Divisibilitat}
% Opcional abans de \input: \def\fmneed{14} (espai mínim abans de cada secció).
% =====================================================================
\makeatletter
\def\fm@logodir{../../logo/pdf}
\providecommand{\fmneed}{10}
\def\fm@cursnom{}

% clau del curs, color i color del text damunt del color
\newcommand{\fm@curs}[3]{\@namedef{fm@color@#1}{#2}\@namedef{fm@text@#1}{#3}}
\fm@curs{1eso}{E6B84A}{navy}
\fm@curs{2eso}{3B82F6}{white}
\fm@curs{3eso}{FF7A50}{white}
\fm@curs{4eso-a}{A78BFA}{white}
\fm@curs{4eso-b}{6D28D9}{white}
\fm@curs{1batx-mi}{9B1C1C}{white}
\fm@curs{1batx-ccss}{EF6B6B}{white}
\fm@curs{2batx-mii}{1E6B47}{white}
\fm@curs{2batx-ccss}{2BB3A3}{white}

\newcommand{\fmcurs}[1]{%
  \@ifundefined{fm@color@#1}{\PackageError{fentmates}{Curs desconegut: #1}{}}{}%
  \edef\fm@tmp{\noexpand\definecolor{curs}{HTML}{\@nameuse{fm@color@#1}}%
               \noexpand\colorlet{curstext}{\@nameuse{fm@text@#1}}}\fm@tmp
  \colorlet{cursfosc}{curs!72!black}%
  \def\fm@logo{\fm@logodir/vertical-#1.pdf}%
  \def\fm@marca{\fm@logodir/marca-#1.pdf}}

% ---------- capçalera de la primera pàgina ----------
% #1 matèria (petit), #2 curs, #3 tipus de document i unitat, #4 títol
\newcommand{\capcalera}[4]{%
  \gdef\fm@cursnom{#2}%
  \par\noindent
  \begin{tcolorbox}[colback=curs,colframe=curs,boxrule=0pt,arc=2.2mm,width=3.5cm,height=2.15cm,
      valign=center,left=3mm,right=2mm,top=1mm,bottom=1mm,nobeforeafter,box align=center]
    \color{curstext}{\footnotesize #1}\par\vspace{1pt}{\fontsize{19}{21}\selectfont\bfseries #2\par}%
  \end{tcolorbox}\hspace{5mm}%
  \begin{minipage}[c]{\dimexpr\linewidth-3.5cm-5mm-2.9cm\relax}
    {\color{cursfosc}\small\bfseries\MakeUppercase{#3}}\par\vspace{2pt}
    {\color{navy}\fontsize{20}{23}\selectfont\bfseries #4\par}
  \end{minipage}\hfill
  \begin{minipage}[c]{2.6cm}\raggedleft\includegraphics[width=2.6cm]{\fm@logo}\end{minipage}\par
  \vspace{2pt}}

% ---------- barres de secció del color del curs ----------
\renewcommand{\seccio}[2]{\par\needspace{\fmneed\baselineskip}\bigskip
  \begin{tcolorbox}[colback=curs,colframe=curs,coltext=curstext,arc=2mm,boxrule=0pt,left=3mm,right=3mm,top=1.2mm,bottom=1.2mm]
  {\large\bfseries #1\quad #2}
  \end{tcolorbox}\vspace{-2pt}}

% ---------- estrelles de dificultat del color del curs (a Batxillerat no se'n posen) ----------
\@ifundefined{estrelles}{}{%
  \renewcommand{\estrelles}[1]{{\color{curs}\ifcase#1\or\faStar\faStar[regular]\faStar[regular]\or\faStar\faStar\faStar[regular]\or\faStar\faStar\faStar\fi}}}

% ---------- capçalera i peu de les pàgines ----------
% primera pàgina: logos institucionals al peu
\fancypagestyle{primera}{\fancyhf{}\renewcommand{\headrulewidth}{0pt}
  \fancyfoot[L]{\raisebox{-0.32\height}{\includegraphics[height=0.78cm]{logo-gva.png}}\hspace{5mm}%
    \raisebox{-0.32\height}{\includegraphics[height=0.78cm]{logo-vermellar.png}}}
  \fancyfoot[R]{\scriptsize\color{muted}\href{https://creativecommons.org/licenses/by-nc-sa/4.0/deed.ca}{\ccbyncsa}\ \ CC BY-NC-SA 4.0 · fentmates.org\qquad\small\thepage}}
% la resta: el símbol a l'esquerra i el curs a la dreta, amb el títol al centre com abans
\fancyhead[L]{\raisebox{-0.3\height}{\includegraphics[height=0.62cm]{\fm@marca}}}
\fancyhead[R]{\small\bfseries\color{cursfosc}\fm@cursnom}
\renewcommand{\headrule}{\color{curs!45}\hrule height \headrulewidth}
\makeatother

%   \fmcurs{1eso}
% i, al començament del document:
%   \capcalera{Matemàtiques}{1r ESO}{Fitxa d'exercicis · Unitat 2}{Divisibilitat}
% Opcional abans de \input: \def\fmneed{14} (espai mínim abans de cada secció).
% =====================================================================
\makeatletter
\def\fm@logodir{../../logo/pdf}
\providecommand{\fmneed}{10}
\def\fm@cursnom{}

% clau del curs, color i color del text damunt del color
\newcommand{\fm@curs}[3]{\@namedef{fm@color@#1}{#2}\@namedef{fm@text@#1}{#3}}
\fm@curs{1eso}{E6B84A}{navy}
\fm@curs{2eso}{3B82F6}{white}
\fm@curs{3eso}{FF7A50}{white}
\fm@curs{4eso-a}{A78BFA}{white}
\fm@curs{4eso-b}{6D28D9}{white}
\fm@curs{1batx-mi}{9B1C1C}{white}
\fm@curs{1batx-ccss}{EF6B6B}{white}
\fm@curs{2batx-mii}{1E6B47}{white}
\fm@curs{2batx-ccss}{2BB3A3}{white}

\newcommand{\fmcurs}[1]{%
  \@ifundefined{fm@color@#1}{\PackageError{fentmates}{Curs desconegut: #1}{}}{}%
  \edef\fm@tmp{\noexpand\definecolor{curs}{HTML}{\@nameuse{fm@color@#1}}%
               \noexpand\colorlet{curstext}{\@nameuse{fm@text@#1}}}\fm@tmp
  \colorlet{cursfosc}{curs!72!black}%
  \def\fm@logo{\fm@logodir/vertical-#1.pdf}%
  \def\fm@marca{\fm@logodir/marca-#1.pdf}}

% ---------- capçalera de la primera pàgina ----------
% #1 matèria (petit), #2 curs, #3 tipus de document i unitat, #4 títol
\newcommand{\capcalera}[4]{%
  \gdef\fm@cursnom{#2}%
  \par\noindent
  \begin{tcolorbox}[colback=curs,colframe=curs,boxrule=0pt,arc=2.2mm,width=3.5cm,height=2.15cm,
      valign=center,left=3mm,right=2mm,top=1mm,bottom=1mm,nobeforeafter,box align=center]
    \color{curstext}{\footnotesize #1}\par\vspace{1pt}{\fontsize{19}{21}\selectfont\bfseries #2\par}%
  \end{tcolorbox}\hspace{5mm}%
  \begin{minipage}[c]{\dimexpr\linewidth-3.5cm-5mm-2.9cm\relax}
    {\color{cursfosc}\small\bfseries\MakeUppercase{#3}}\par\vspace{2pt}
    {\color{navy}\fontsize{20}{23}\selectfont\bfseries #4\par}
  \end{minipage}\hfill
  \begin{minipage}[c]{2.6cm}\raggedleft\includegraphics[width=2.6cm]{\fm@logo}\end{minipage}\par
  \vspace{2pt}}

% ---------- barres de secció del color del curs ----------
\renewcommand{\seccio}[2]{\par\needspace{\fmneed\baselineskip}\bigskip
  \begin{tcolorbox}[colback=curs,colframe=curs,coltext=curstext,arc=2mm,boxrule=0pt,left=3mm,right=3mm,top=1.2mm,bottom=1.2mm]
  {\large\bfseries #1\quad #2}
  \end{tcolorbox}\vspace{-2pt}}

% ---------- estrelles de dificultat del color del curs (a Batxillerat no se'n posen) ----------
\@ifundefined{estrelles}{}{%
  \renewcommand{\estrelles}[1]{{\color{curs}\ifcase#1\or\faStar\faStar[regular]\faStar[regular]\or\faStar\faStar\faStar[regular]\or\faStar\faStar\faStar\fi}}}

% ---------- capçalera i peu de les pàgines ----------
% primera pàgina: logos institucionals al peu
\fancypagestyle{primera}{\fancyhf{}\renewcommand{\headrulewidth}{0pt}
  \fancyfoot[L]{\raisebox{-0.32\height}{\includegraphics[height=0.78cm]{logo-gva.png}}\hspace{5mm}%
    \raisebox{-0.32\height}{\includegraphics[height=0.78cm]{logo-vermellar.png}}}
  \fancyfoot[R]{\scriptsize\color{muted}\href{https://creativecommons.org/licenses/by-nc-sa/4.0/deed.ca}{\ccbyncsa}\ \ CC BY-NC-SA 4.0 · fentmates.org\qquad\small\thepage}}
% la resta: el símbol a l'esquerra i el curs a la dreta, amb el títol al centre com abans
\fancyhead[L]{\raisebox{-0.3\height}{\includegraphics[height=0.62cm]{\fm@marca}}}
\fancyhead[R]{\small\bfseries\color{cursfosc}\fm@cursnom}
\renewcommand{\headrule}{\color{curs!45}\hrule height \headrulewidth}
\makeatother

\fmcurs{4eso-a}

\begin{document}
\thispagestyle{primera}

\capcalera{Matemáticas A}{4\textsuperscript{o} ESO}{Ficha de ejercicios · Unidad 4}{\raggedright Ecuaciones}
\vspace{10pt}
Nombre y apellidos: \rule{6.6cm}{0.4pt}\hfill Grupo: \rule{1.6cm}{0.4pt}\hfill Fecha: \rule{2.2cm}{0.4pt}
\vspace{8pt}

\begin{tcolorbox}[colback=soft,colframe=soft,arc=2mm,left=3mm,right=3mm,top=1.5mm,bottom=1.5mm]
\small Ejercicios graduados por niveles de dificultad:\quad \estrelles{1}\ básico \qquad \estrelles{2}\ medio \qquad \estrelles{3}\ avanzado.\\[2pt] Resuélvelos en el cuaderno mostrando todos los pasos. Puedes comprobar los resultados con el resolutor paso a paso de los apuntes interactivos.
\end{tcolorbox}

% =====================================================
\seccio{1}{Ecuaciones de primer grado}

\exf{1} Resuelve:\begin{tasks}(3)
  \task $3x-7=2x+5$
  \task $5x+4=2x-11$
  \task $2(x-3)=x+4$
  \task $4-(x+2)=3x-6$
  \task $3(2x-1)-2(x+4)=1$
  \task $7-2(3-x)=5(x-1)$
\end{tasks}\recorda{Primero quita los paréntesis; después pasa las $x$ a un miembro y los números al otro.}

\ex{2} Resuelve e indica si alguna ecuación no tiene solución o tiene infinitas:\begin{tasks}(2)
  \task $5(x-2)-3(2x+1)=2(x-8)$
  \task $2[x-3(x+1)]=4-x$
  \task $x-2[x-(3-x)]=7$
  \task $3(x+2)-(x-4)=2(x+5)$
  \task $4(x-1)-2(x+3)=2(x-6)$
  \task $6-[2x-(x+5)]=3(x-1)$
\end{tasks}\recorda{Cuidado con el signo menos delante de un paréntesis o de un corchete: $-(x-3)=-x+3$. Empieza por los paréntesis de dentro.}

\ex{2} Resuelve (puedes multiplicar por 10 para eliminar los decimales):\begin{tasks}(2)
  \task $0,2(x+5)-0,5x=0,1(x-6)$
  \task $1,5(x-2)=0,5x+3$
  \task $0,3x-1,2=0,1(x+4)$
\end{tasks}

\ex{2} Comprueba, sin resolverlas, si $x=-2$ es solución de cada ecuación:\begin{tasks}(2)
  \task $3(x+4)=x+10$
  \task $x^{2}-3x=2(3-x)$
  \task $\dfrac{x+1}{3}=\dfrac{x-4}{2}$
  \task $4-3x=5(x+4)$
\end{tasks}

\ex{3} Escribe una ecuación de primer grado con paréntesis que tenga por solución $x=-3$, una que no tenga solución y una que tenga infinitas.


% =====================================================
\seccio{2}{Ecuaciones de primer grado con denominadores}

\exf{1} Resuelve:\begin{tasks}[after-item-skip=12pt](2)
  \task $\dfrac{x}{2}+\dfrac{x}{3}=10$
  \task $x-\dfrac{x}{4}=9$
  \task $\dfrac{x+1}{3}=\dfrac{x-1}{2}$
  \task $\dfrac{x}{3}-\dfrac{1}{2}=\dfrac{x}{4}$
\end{tasks}

\ex{2} Resuelve:\begin{tasks}[after-item-skip=14pt](2)
  \task $\dfrac{x-2}{3}-\dfrac{x+1}{4}=1$
  \task $\dfrac{3x-1}{2}-\dfrac{x+3}{5}=x+1$
  \task $x-\dfrac{x-2}{6}=\dfrac{2x+1}{3}+\dfrac{1}{2}$
  \task $\dfrac{5x+1}{4}-\dfrac{x-1}{2}=\dfrac{x+5}{4}$
\end{tasks}\recorda{Multiplica todos los términos por el m.c.m. de los denominadores. Si el numerador tiene varios términos, ponlo entre paréntesis: $-\dfrac{x-1}{2}\cdot4=-2(x-1)$.}

\ex{3} Resuelve e indica si alguna ecuación no tiene solución o tiene infinitas:\begin{tasks}[after-item-skip=14pt](2)
  \task $\dfrac{2(x-1)}{3}-\dfrac{3x+1}{4}=\dfrac{x-5}{6}-1$
  \task $\dfrac{x+1}{2}-\dfrac{2x-3}{3}=\dfrac{x+3}{6}-\dfrac{x-1}{3}$
  \task $\dfrac{3(x-2)}{4}-\dfrac{x+1}{2}=\dfrac{x-8}{4}$
  \task $\dfrac{x}{2}-\dfrac{x-\dfrac{x+1}{3}}{4}=1$
\end{tasks}

\ex{3} Resuelve:\begin{tasks}[after-item-skip=14pt](1)
  \task $\dfrac{x+4}{5}-\dfrac{x-2}{10}=\dfrac{x+1}{2}-\dfrac{2x-1}{4}$
  \task $1-\dfrac{x-3}{6}=\dfrac{2-x}{4}+\dfrac{x+5}{12}$
\end{tasks}


% =====================================================
\seccio{3}{Problemas de primer grado}

\exf{1} La suma de tres números consecutivos es 147. ¿Cuáles son?

\ex{1} Un abrigo rebajado un 20\,\% cuesta 68 €. ¿Cuánto costaba antes de la rebaja?

\ex{2} Laura tiene 14 años y su madre, 41. ¿Hace cuántos años la edad de la madre era el cuádruple de la de Laura?

\ex{2} Un triángulo isósceles tiene 50 cm de perímetro. El lado desigual mide 4 cm menos que cada uno de los lados iguales. ¿Cuánto mide cada lado?

\ex{2} De un depósito de agua se gasta $\frac13$ el lunes y el martes, $\frac25$ de lo que quedaba. Todavía quedan 240 L. ¿Qué capacidad tiene el depósito?

\ex{2} Mezclamos 30 kg de café de 8 €/kg con café de 12 €/kg y obtenemos una mezcla que sale a 9,50 €/kg. ¿Cuántos kilos del segundo café hemos puesto?

\ex{2} Un padre reparte 1\,400 € entre tres hijos. El segundo recibe el doble que el primero, y el tercero, 200 € más que el segundo. ¿Cuánto recibe cada uno?

\ex{3} Un coche sale de una ciudad a 90 km/h. Una hora y media después sale otro por el mismo camino a 120 km/h. ¿Cuánto tarda el segundo en alcanzar al primero? ¿A qué distancia de la ciudad lo alcanza?

\ex{3} Un grifo llena un depósito en 6 horas y otro, en 3 horas. ¿Cuánto tardan en llenarlo si los abrimos los dos a la vez?


% =====================================================
\seccio{4}{Ecuaciones de segundo grado incompletas}

\exf{1} Resuelve las ecuaciones incompletas:\begin{tasks}(3)
  \task $x^{2}-49=0$
  \task $3x^{2}-27=0$
  \task $x^{2}+16=0$
  \task $5x^{2}-20x=0$
  \task $2x^{2}=7x$
  \task $4x^{2}-25=0$
  \task $x^{2}=5$
  \task $-x^{2}+6x=0$
  \task $9x^{2}+4=0$
\end{tasks}\recorda{$ax^2+c=0\Rightarrow x^2=-\frac ca$ (dos soluciones opuestas o ninguna). \ $ax^2+bx=0\Rightarrow x(ax+b)=0\Rightarrow x=0,\ x=-\frac ba$.}

\ex{2} Opera y resuelve (todas acaban en una ecuación incompleta):\begin{tasks}[after-item-skip=10pt](2)
  \task $(x+3)^{2}=6x+18$
  \task $x(x-4)=2x(x+1)$
  \task $(2x-1)(2x+1)=15$
  \task $(x-2)^{2}+(x+2)^{2}=26$
  \task $\dfrac{x^{2}}{3}-\dfrac{x}{2}=0$
  \task $(x+5)(x-5)=\dfrac{x^{2}}{2}-5$
\end{tasks}

\ex{3} ¿Para qué valores de $c$ tiene soluciones la ecuación $x^2+c=0$? Y la ecuación $2x^2+bx=0$, ¿cuántas tiene según el valor de $b$?


% =====================================================
\seccio{5}{Ecuaciones de segundo grado completas}

\exf{1} Resuelve:\begin{tasks}(3)
  \task $x^{2}-8x+15=0$
  \task $x^{2}+3x-10=0$
  \task $2x^{2}-5x+2=0$
  \task $x^{2}-10x+25=0$
  \task $3x^{2}+x+2=0$
  \task $6x^{2}+x-2=0$
\end{tasks}\recorda{$ax^2+bx+c=0\Rightarrow x=\dfrac{-b\pm\sqrt{b^2-4ac}}{2a}$.}

\ex{2} Resuelve y deja las soluciones con radicales:\begin{tasks}(3)
  \task $x^{2}-4x+1=0$
  \task $x^{2}+2x-4=0$
  \task $2x^{2}-6x+3=0$
\end{tasks}

\ex{2} Elimina los decimales o los denominadores y resuelve:\begin{tasks}[after-item-skip=10pt](3)
  \task $0,5x^{2}-1,5x+1=0$
  \task $\dfrac{x^{2}}{4}-x+1=0$
  \task $0,2x^{2}+0,6x=2$
\end{tasks}

\ex{2} Sin resolverlas, calcula el discriminante e indica cuántas soluciones tiene cada ecuación:\begin{tasks}(3)
  \task $2x^{2}-3x+5=0$
  \task $x^{2}+12x+36=0$
  \task $3x^{2}-x-1=0$
\end{tasks}

\ex{3} Halla el valor de $k$ para que la ecuación $x^2+kx+9=0$ tenga una solución doble, y calcúlala.

\ex{3} Escribe una ecuación de segundo grado con coeficientes enteros que tenga por soluciones $x=2$ y $x=-\frac32$.


% =====================================================
\seccio{6}{Segundo grado con paréntesis y denominadores}

\exf{2} Opera hasta obtener la forma $ax^2+bx+c=0$ y resuelve:\begin{tasks}[after-item-skip=6pt](2)
  \task $(x-3)^{2}=2x-6$
  \task $(2x+1)(x-2)=x^{2}-4$
  \task $x(x+2)-3(x-1)=9$
  \task $(x+1)^{2}-(x-2)(x+2)=x^{2}-3$
\end{tasks}

\ex{2} Elimina los denominadores y resuelve:\begin{tasks}[after-item-skip=14pt](2)
  \task $\dfrac{x^{2}}{2}-\dfrac{x}{3}=\dfrac{x}{6}+1$
  \task $\dfrac{x(x+1)}{3}-\dfrac{x+2}{2}=0$
  \task $\dfrac{(x+3)(x-3)}{4}-\dfrac{(x-1)^{2}}{2}=-2$
  \task $\dfrac{(x-2)^{2}}{3}-\dfrac{x+1}{2}=\dfrac{x-11}{6}$
\end{tasks}\recorda{Primero desarrolla los productos y las identidades notables, y después multiplica por el m.c.m. ¡Cuidado con el signo menos delante de una fracción!}

\ex{3} Resuelve:\begin{tasks}[after-item-skip=14pt](1)
  \task $\dfrac{(x+1)^{2}}{2}-\dfrac{(x-1)^{2}}{3}=\dfrac{x+11}{6}+x$
  \task $\dfrac{(2x-1)^{2}}{3}-\dfrac{(x+1)(x-1)}{2}=\dfrac{x+7}{6}$
  \task $\dfrac{(x-1)^{2}}{2}-\dfrac{(x+1)(x-3)}{3}=\dfrac{x+9}{6}$
  \task $2(x-1)^{2}-3(x+1)(x-2)=x(x-5)-4$
\end{tasks}


% =====================================================
\seccio{7}{Problemas de segundo grado}

\exf{2} Si a un número natural le sumamos su cuadrado obtenemos 110. ¿Cuál es el número?

\ex{2} El producto de las edades de dos hermanos es 168, y uno tiene 2 años más que el otro. ¿Cuántos años tienen?

\ex{2} Un rectángulo tiene 34 m de perímetro y 60 m² de área. ¿Qué dimensiones tiene?

\ex{2} La diagonal de un rectángulo mide 13 cm y un lado mide 7 cm más que el otro. Calcula sus lados y el área.

\ex{2} Si un lado de un jardín cuadrado lo alargamos 3 m y el otro lo acortamos 2 m, obtenemos un rectángulo de 66 m². ¿Cuánto mide el lado del jardín?

\ex{3} Una foto de $20\times15$ cm tiene un marco de la misma anchura por todos los lados. El área de la foto con el marco es de 594 cm². ¿Qué anchura tiene el marco?

\ex{3} Un grupo de amigos tiene que pagar 120 € a partes iguales. Si fueran 2 amigos más, cada uno pagaría 5 € menos. ¿Cuántos amigos son?

\ex{3} Los lados de un triángulo rectángulo son tres números naturales consecutivos. ¿Cuánto miden?


% =====================================================
\seccio{8}{Ecuaciones factorizadas}

\exf{1} Resuelve sin multiplicar los factores:\begin{tasks}(3)
  \task $(x-3)(x+5)=0$
  \task $x(2x-7)=0$
  \task $(x+1)(x-2)(3x+4)=0$
  \task $5(x-6)^{2}=0$
  \task $x(4x-1)(x+4)=0$
  \task $(x^{2}-9)(x-1)=0$
\end{tasks}\recorda{Un producto vale cero si alguno de los factores vale cero: iguala cada factor a 0.}

\ex{2} Resuelve:\begin{tasks}(2)
  \task $(x^{2}-5x+6)(x+4)=0$
  \task $(x^{2}+4)(x^{2}-16)=0$
  \task $x^{2}(x^{2}-3x-10)=0$
  \task $(2x-1)^{2}(x^{2}-2)=0$
\end{tasks}

\ex{2} Factoriza sacando factor común (y, si hace falta, con una identidad notable) y resuelve:\begin{tasks}(3)
  \task $x^{3}-9x=0$
  \task $2x^{3}-8x^{2}=0$
  \task $x^{3}-10x^{2}+25x=0$
  \task $x^{4}-x^{2}=0$
  \task $x^{3}+x^{2}-6x=0$
  \task $3x^{3}-12x=0$
\end{tasks}

\ex{3} Resuelve (cuidado: no están igualadas a cero):\begin{tasks}(2)
  \task $(x-1)(x+2)=x-1$
  \task $x(x-3)=2(x-3)$
\end{tasks}


% =====================================================
\seccio{9}{Ecuaciones racionales}

\exf{1} Resuelve:\begin{tasks}[after-item-skip=12pt](3)
  \task $\dfrac{5}{x}-1=4$
  \task $\dfrac{x}{x-2}=3$
  \task $\dfrac{x+1}{x-1}=3$
  \task $\dfrac{2}{x+1}=\dfrac{3}{x-1}$
  \task $\dfrac{x}{4}=\dfrac{9}{x}$
  \task $\dfrac{x+2}{x}=\dfrac{x+8}{x+2}$
\end{tasks}\recorda{En una proporción $\dfrac AB=\dfrac CD$ se puede multiplicar en cruz: $A\cdot D=B\cdot C$.}

\ex{2} Resuelve:\begin{tasks}[after-item-skip=12pt](2)
  \task $x+\dfrac{8}{x}=6$
  \task $\dfrac{1}{x}+\dfrac{1}{x+3}=\dfrac{1}{2}$
  \task $\dfrac{x}{x-1}+\dfrac{2}{x}=3$
  \task $\dfrac{3}{x-1}+\dfrac{2}{x+1}=2$
\end{tasks}

\ex{3} Resuelve y comprueba las soluciones:\begin{tasks}[after-item-skip=14pt](2)
  \task $\dfrac{x}{x-1}-\dfrac{2}{x+1}=\dfrac{4}{x^{2}-1}$
  \task $\dfrac{x}{x-3}+1=\dfrac{3}{x-3}$
  \task $\dfrac{6}{x^{2}-4}-\dfrac{1}{x-2}=\dfrac{1}{x+2}$
  \task $\dfrac{x+1}{x}-\dfrac{x}{x+1}=\dfrac{5}{6}$
\end{tasks}\recorda{Antes de dar las soluciones, descarta las que anulan algún denominador.}

\ex{3} La suma de un número y su inverso es $\frac{10}{3}$. ¿Cuál es el número?


% =====================================================
\seccio{10}{Ecuaciones con radicales}

\exf{1} Resuelve:\begin{tasks}(3)
  \task $\sqrt{x+3}=4$
  \task $\sqrt{2x-1}=3$
  \task $\sqrt{x^{2}+9}=5$
  \task $\sqrt{3x+1}=-2$
  \task $2\sqrt{x}=6$
  \task $\sqrt{5x-1}=2$
\end{tasks}

\ex{2} Resuelve y comprueba las soluciones:\begin{tasks}(3)
  \task $\sqrt{x+7}=x+1$
  \task $x-\sqrt{x+2}=4$
  \task $\sqrt{2x+1}+1=x$
  \task $\sqrt{3x-2}=x$
  \task $x+\sqrt{x-1}=3$
  \task $2\sqrt{x+1}=x+1$
\end{tasks}\recorda{Aísla el radical, eleva los dos miembros al cuadrado y comprueba todas las soluciones en la ecuación inicial.}

\ex{2} Resuelve (dos radicales: pon uno en cada miembro y eleva al cuadrado):\begin{tasks}(3)
  \task $\sqrt{2x+3}=\sqrt{x+5}$
  \task $\sqrt{5x-4}-\sqrt{x+8}=0$
  \task $3\sqrt{x-1}=\sqrt{x+7}$
  \task $\sqrt{x^{2}-7}=\sqrt{x+5}$
  \task $\sqrt{x^{2}+3x}=2\sqrt{x}$
  \task $\sqrt{x^{2}-3}=\sqrt{2x}$
\end{tasks}

\ex{3} Resuelve y comprueba las soluciones:\begin{tasks}(2)
  \task $x-\sqrt{25-x^{2}}=1$
  \task $\sqrt{x^{2}+5}+1=2x$
\end{tasks}

\ex{3} Si a un número le sumamos 2 y hacemos su raíz cuadrada, obtenemos el mismo número menos 4. ¿Cuál es el número?

\vfill
\begin{center}\color{muted}\small\faCheckCircle[regular]\ \ ¡Ánimo! Comprueba tus resultados con las herramientas interactivas de los apuntes de la unidad.\end{center}

\end{document}
